\section{The Observational Foundation: Longer, Denser, and Harmonised Monitoring}

Across the drainage literature, the field's most secure statements are
about the limits of its own evidence. Headline results are drawn
disproportionately from short records, single sites, and narrow
agro-ecological windows, and their authors caution explicitly against
premature generalisation. The recurring signal is therefore not a settled
conclusion but a precondition: records long enough, dense enough,
spatially replicated enough, and consistent enough to separate a
management signal from the natural variability of weather, hydrology, and
slow soil change. \textcite{Rogger2017} pitch this at the level of the
whole discipline, calling for complex-systems thinking that links
processes across time scales, long-term experiments on
physical--chemical--biological process interactions, attention to
connectivity and patterns across spatial scales, typologies and
meta-analyses that generalise findings, and multidisciplinary consortia
sharing data and terminology under known boundary conditions. The sheer
frequency of blunt appeals---that ``more research is needed''
\parencite{SHALIT1995}, that ``the complexity of this subject requires more
research'' \parencite{Malone2014}, that ``further research in this area is
needed'' \parencite{Woo2023}, and that ``further evaluations are warranted
before any concrete conclusions could be drawn''
\parencite{GUPTA1993}---marks genuine gaps in evidence rather than
rhetorical closings. Because every later domain in this review inherits
these limits, the observational foundation is treated here as a frontier
in its own right.

\subsection{Extending the temporal record}

The single most widely voiced need is for records long enough to
disentangle treatment effects from climatic noise and from slow,
directional change. \textcite{Clifford2023} emphasise truly long-term
experiments, and \textcite{BUHLER1993} caution that a lag of several years
may precede any objective evaluation of an altered practice. Minimum
horizons are proposed---at least five years, with continued observation as
systems are extended in length or depth over time \parencite{Ali2013}, and
lengthening national time series expected to sharpen analysis as records
grow \parencite{Ali2001}. Many studies attach conclusions to explicit
requirements for further years: to quantify a denitrification wall's
effectiveness and its long-term cost per unit nitrate removed against
natural variability \parencite{Jaynes2008}, to establish how long treatment
benefits persist \parencite{Dalhem2019,Osterholm2005}, and to test the
replicability of results in time and confirm that observed differences
endure under comparable conditions
\parencite{ROLLIN1987,Reinhart2020,Rijal2012,GiménezValero2025}, all
echoing the disciplinary case for long-term monitoring at multiple scales
\parencite{SPALING1995}.

Nutrient and hydrological studies converge on the same horizon.
Year-round measurement of flow and concentration for at least several years
on new sites is deemed necessary to set management priorities
\parencite{Kladivko2001}, and long-term records are singled out to close
water and nutrient balances and study legacy pools \parencite{Sharpley2019}
and to disentangle land-management from climate-change effects, with the
corollary that policy evaluation demands patience
\parencite{Bauwe2020,Kelley2017,Gillingham2000,Lemke2011}. Smoothing
year-to-year variability requires evaluating both phases of a rotation
\parencite{Randall2003,Randall2005}, assessing responses on an annualised,
multi-year basis \parencite{Waring2024,Pittelkow2017}, and freeing
drainage-water-management yield effects from confounding
\parencite{Cooke2012}; only long-term records across a range of climatic
conditions, gathering plant, soil, and water measurements together, yield
appropriate guidelines \parencite{RANDALL1995}. Data spanning at least a
normal, a wet, and a dry year are needed to resolve losses and transport
routes \parencite{vanderSalm2012}, and as long a record as possible is
required for trend estimation because water quality may lag the
management-driven change \parencite{Ulén2007}. Long records are equally the
only way to detect slow processes---soil organic-carbon recovery
\parencite{Kumar2014}, discharge-driven organic-matter export
\parencite{Manninen2023}, and whether drainage-induced soil evolution
continues or stabilises over decades \parencite{Frison2009,Frankl2023}.
Record length alone is insufficient: robust seasonal conclusions require
year-round coverage, especially in the non-growing season
\parencite{Klaiber2020,VanEsbroeck2016}, high-frequency records without
gaps and instruments checked as often as possible \parencite{Hodaj2017},
and monitoring timed to infrastructure interventions such as drain repair
\parencite{Wesström2015}.

\subsection{Broadening spatial and geographic scope}

A second frontier is testing whether site-specific findings hold across
other soils, climates, crops, and management regimes. Authors repeatedly
defer broad adoption until similar data are collected from further sites and
years \parencite{Menegaz2025,Zimmerman2017b,Terrado2015}, from additional
agro-ecological zones \parencite{Ali2001}, from other settings for
practices such as the two-stage ditch \parencite{Hodaj2017}, and from
adverse or compacted conditions where a technology's merit is untested
\parencite{Bigelow2001}. The most explicit prescription is that
controlled-drainage and subirrigation studies be conducted outside the U.S.
Midwest and Canada across diverse climates, crops, and soils, and---
crucially for generalisation---that site characterisation (slope,
particle-size distribution by horizon, hydraulic conductivity, drainage
class, and moisture) be reported with measures of variability
\parencite{Singh2022b}. Repetition at larger scales and on contrasting soils
is stressed as a key modulating factor
\parencite{Bilotta2008,Goeken2015}, including for process studies such as
earthworm activity on soils of differing texture and rainfall
\parencite{Nuutinen2001,Nuutinen2011}. Climate-focused transferability
demands high-resolution scenarios and multi-model uncertainty analysis at
additional sites \parencite{Rasmussen2018,Jeantet2022,Jin2008}, and testing
across regions, soils, and climates recurs throughout as the precondition
for localised design guidance
\parencite{Singh2020b,Malota2024,Salazar2008,Holbak2021}. At the level of
synthesis, robust conclusions will require geographically diverse and
methodologically harmonised studies so that meta-analyses can disentangle
management effects \parencite{Wu2025}.

\subsection{Monitoring infrastructure, instrumentation, and data tools}

A distinct set of needs concerns the apparatus that makes long-term data
collection feasible and reliable. \textcite{Kanwar1999} describe how
automatic samplers collect discrete alongside continuous samples, how data
loggers can be extended to soil temperature and moisture, and how relays or
pressure switches could detect sump-pump operation, advising that
equipment-to-logger distances be minimised, data downloaded automatically to
reduce loss, and annual maintenance---cleaning check valves and calibrating
flowmeters---be performed. Diagnostic methods are themselves research
targets: reducing pressure variations by doubling equipment and injecting
smoke from both ends of a drain section could enable work with longer drain
lines, though transport pathways and observed clustering remain unresolved
\parencite{Petersen2012}, and selected pipes, envelopes, perforations, and
surrounding soil should be observed with permanent monitoring of the
sampling process to confirm favourable results are maintained
\parencite{MartínezBeltrán2007}. Closing the water budget requires
high-frequency discharge, runoff, and seepage measurement, process
measurements of microbial activity, denitrification, and redox, and a
publicly available dataset \parencite{Frankenberger2024}. Software and data
infrastructure are an increasingly prominent part of this frontier:
readers are directed to evolving data products such as the National
Elevation Dataset and USGS resources \parencite{Shedekar2018}, and a series
of drainage-network configuration and sizing tools describe themselves as
under continual development and openly invite testing with additional
datasets \parencite{Falasy2025,Falasy2025b,Falasy2024}.

\subsection{Tracers, isotopes, and source apportionment}

Measurement and tracing techniques are repeatedly identified as the means
to resolve pathways that bulk flow-and-concentration records leave
ambiguous. \textcite{Williams2022} call for higher-frequency (sub-daily to
daily) measurements, further work on groundwater dynamics in dissolved-P
transport, and more comprehensive source-apportionment methods.
Water- and nitrate-isotope approaches are advanced to quantify processes
directly: coupling water and nitrate oxygen isotopes and quantifying the
$\delta^{18}$O of nitrification-derived nitrate to constrain evaporative
enrichment \parencite{Hu2024}, monitoring nitrate concentrations and
isotopic compositions with multiple tracers across topsoil, drains, and
ditches to characterise the soil contribution to river nitrate
\parencite{Briand2017}, and boron isotopes to strengthen contaminant
tracing around feedlots and fields \parencite{Komor1997}. High-frequency
tracer measurements are urged to quantify water ages and represent both
celerities and velocities \parencite{Williams2021}, alongside event
sampling and bulk-soil-water isotopes to constrain young-water fractions
\parencite{DimitrovaPetrova2020}, high-resolution specific conductance for
end-member mixing \parencite{Kronholm2015,Schilling2019}, and
conservative-tracer or water-isotope studies to partition quick flow and
baseflow \parencite{Frederiksen2023}. Explicit pathway partitioning---
reporting surface-inlet presence and splitting water and nutrient loads by
pathway via before-after control-impact designs, hydrograph separation, or
tracers \parencite{Flores2021}, evaluating overland-flow versus tile-flow
contributions \parencite{Kokulan2019}, and verifying groundwater flow paths
and plant water sources with stable isotopes
\parencite{Islam2022,YliHalla2022}---recurs throughout. The role of
groundwater outflow in nutrient loading is flagged as under-quantified and
in need of dedicated measurement and modelling
\parencite{Turunen2013,Jokinen2024,Barkle2021}, and finer temporal analyses
of the flow--concentration relationship, including detailed storm-hydrograph
separation, are needed to test event-scale hypotheses
\parencite{Lagzdins2015,Danalatos2022,Kladivko2004}.

\subsection{Standardisation, harmonisation, and comparability}

For accumulating data to support synthesis, methods and reporting must be
harmonised. The value of future meta-analyses is tied directly to
methodologically harmonised, comprehensive, long-term monitoring of the full
suite of environmental impacts \parencite{Wu2025}. \textcite{Zhang2018} set
out the most detailed programme: improving the quality and comparability of
existing indicators through harmonised data collection, monitoring,
interpretation, and reporting; developing new indicators for emerging
concerns while ensuring scientific credibility, sensitivity to change, and
ease of interpretation; establishing standardised evaluation procedures
including thresholds and reference points; recognising data gaps and
improving upscaling and uncertainty assessment; and incorporating
earth-observation data. These priorities complement the call for typologies,
meta-analyses, and shared data and terminology across research consortia
\parencite{Rogger2017}. At the study level, meta-analytic comparability
requires complete hydrologic description of each plot, standardised
analytical methods, adequate monitoring periods, and full reporting of
concentrations with descriptive statistics \parencite{Ghirardini2019}, and
soil biological and soil-health indicators must be sampled at consistent
times under similar moisture and temperature conditions to be comparable
across sites \parencite{Sherbine2023,Li2023b}.

\subsection{Experimental design for detectable effects}

Longer, standardised records are useful only if the design can resolve the
effect sought. \textcite{Schilling2013} specify the ingredients for
paired-basin detection: establishing a suitable correlation among paired
basins across climates, restricting paired sampling to the March--July
high-flow period, delaying post-treatment monitoring to enhance the chance
of detecting change, and adopting practices capable of achieving at least
the minimum detectable change. Effects should be assessed simultaneously at
multiple scales, from field to watershed \parencite{Sunohara2016}, paired
with within-field and within-plot studies to elucidate the mechanistic
processes of export and additional field-scale studies within high-export
sub-watersheds \parencite{Irvine2019}. Larger, more representative trial
areas are sought to test effective-yield formulas \parencite{Xin2016c} and
to refine tracer-based ($^{15}$N) loss observation \parencite{Xiong2025},
and mesocosms are proposed to quantify nutrient-loss responses to
management change under controlled conditions \parencite{Loper2024}.
Because field-scale drainage trials are typically unreplicated, several
authors treat the analysis itself as an open question: whether beta
distributions, time series, or new methods best analyse unreplicated paired
experiments \parencite{Pagliari2021}, and whether directed hypothesis tests
should be favoured over replicate spatial analyses \parencite{Andrus2014}.
Where replication is feasible, replicated field experiments, matched
instrumentation, and improved sampling designs are recommended to support
statistical inference
\parencite{Okuda2020,Wamono2016b,Davis2019,Jacquemin2020}. Together these
design refinements define the bridge between a lengthening, harmonised
observational record and the detectable, transferable conclusions the rest
of the field depends upon.
